\section{Podstawy teoretyczne zastosowanych metod}
\label{sec:podstawy}

W tym rozdziale opisano --- w sposób możliwie przystępny, bez zagłębiania się w
formalizm matematyczny --- dwie rodziny metod, które porównujemy: las losowy
oraz sieci głębokie. Celem jest zrozumienie, \emph{jak} działają i \emph{czym}
się różnią, bo to te różnice decydują o ich przydatności w ochronie danych w
chmurze.

\subsection{Las losowy}

Punktem wyjścia jest \emph{drzewo decyzyjne}. Można je sobie wyobrazić jako
listę pytań typu ,,czy liczba pakietów w połączeniu przekracza pewien próg?''.
Odpowiadając na kolejne pytania, schodzimy w dół drzewa aż do liścia, który
mówi: ,,to ruch normalny'' albo ,,to atak''. Pojedyncze drzewo jest jednak
kapryśne --- drobna zmiana danych potrafi je całkowicie przebudować, przez co
łatwo się ,,przeucza'' i słabo radzi sobie z nowymi przypadkami.

\emph{Las losowy}~\cite{breiman2001random} rozwiązuje ten problem prostą, ale
skuteczną sztuczką: zamiast jednego drzewa buduje ich setki, a każde uczy na
nieco innym, losowym wycinku danych i cech. Decyzja końcowa zapada przez
\emph{głosowanie} --- jeśli większość drzew uznaje połączenie za atak, takim
jest klasyfikowane. Dzięki temu, że drzewa popełniają różne błędy, ich
uśrednienie daje wynik znacznie stabilniejszy niż każde drzewo z osobna.

Las losowy ma kilka cech, które czynią go popularnym w cyberbezpieczeństwie:
nie wymaga skalowania danych, dobrze radzi sobie ze złożonymi zależnościami
między cechami, szybko się uczy i --- co istotne dla analityka --- potrafi
wskazać, \textbf{które cechy} były najważniejsze przy podejmowaniu decyzji.

\subsection{Sieci głębokie}

\emph{Sieć neuronowa} to model luźno inspirowany działaniem mózgu. Składa się z
warstw ,,neuronów'', z których każdy łączy informacje z poprzedniej warstwy i
przekazuje wynik dalej. Określenie \emph{głęboka} oznacza po prostu, że warstw
jest wiele --- dzięki temu sieć stopniowo buduje coraz bardziej złożone wzorce:
od prostych zależności po abstrakcyjne ,,sygnatury'' złośliwej aktywności.
Sieć uczy się, wielokrotnie oglądając przykłady i samodzielnie korygując swoje
,,wewnętrzne pokrętła'' tak, aby coraz rzadziej się mylić.

W pracy porównujemy trzy architektury sieci:

\paragraph{MLP (perceptron wielowarstwowy).} Najprostsza i najbardziej klasyczna
sieć w pełni połączona. Traktuje dane jak zwykłą tabelę liczb i jest naturalnym
wyborem dla danych opisujących połączenia sieciowe.

\paragraph{CNN 1D (sieć konwolucyjna).} Sieć, która ,,przesuwa okienko'' po
danych i wyszukuje powtarzające się, lokalne wzorce~\cite{lecun1998gradient}.
Sprawdza się świetnie w analizie obrazów czy sygnałów; tutaj sprawdzamy, czy
podobne podejście pomaga w danych o ruchu sieciowym.

\paragraph{LSTM (sieć rekurencyjna).} Sieć z ,,pamięcią'', zaprojektowana do
danych o charakterze sekwencji~\cite{hochreiter1997long}. Potrafi uwzględniać
kolejność i zależności między kolejnymi elementami, dlatego włączamy ją do
porównania, by sprawdzić, czy modelowanie ,,sekwencyjności'' cech przynosi
jakąkolwiek korzyść.

\subsection{Miary oceny skuteczności}

Sama ,,dokładność'' (odsetek poprawnych decyzji) bywa myląca, gdy ataki są
rzadkie --- model, który wszystko uzna za ruch normalny, może mieć wysoką
dokładność, a jednocześnie nie wykryć żadnego włamania. Dlatego patrzymy na
kilka uzupełniających się miar, opisanych tu zwykłym językiem:

\begin{itemize}
  \item \textbf{Precyzja} --- spośród połączeń oznaczonych jako atak, jaki
        odsetek faktycznie był atakiem. Niska precyzja oznacza wiele
        \emph{fałszywych alarmów}, które męczą zespół bezpieczeństwa.
  \item \textbf{Czułość} (ang.\ \emph{recall}) --- spośród wszystkich
        prawdziwych ataków, ile udało się wykryć. Niska czułość oznacza
        \emph{przeoczone włamania} --- najgroźniejszy rodzaj błędu.
  \item \textbf{F1} --- jedna liczba godząca precyzję i czułość; wygodna do
        porównań.
  \item \textbf{ROC-AUC} i \textbf{PR-AUC} --- miary opisujące jakość modelu
        niezależnie od ustawionego ,,progu czułości''; im bliżej $1{,}0$, tym
        lepiej.
\end{itemize}

W kontekście ochrony danych w chmurze szczególnie zależy nam na \emph{wysokiej
czułości} (żeby nie przeoczyć naruszenia) przy rozsądnej precyzji (żeby nie
zasypać analityków fałszywymi alarmami).

\subsection{Przegląd literatury}

Wykrywanie włamań metodami uczenia maszynowego to dobrze zbadany temat.
Najczęściej wykorzystuje się publiczne zbiory danych, takie jak
NSL-KDD~\cite{tavallaee2009detailed}, UNSW-NB15~\cite{moustafa2015unsw} czy
CIC-IDS2017/CSE-CIC-IDS2018 przygotowane przez Canadian Institute for
Cybersecurity~\cite{sharafaldin2018toward}. Co ciekawe, w wielu badaniach metody
oparte na drzewach (jak las losowy) okazują się równie dobre lub lepsze od sieci
głębokich, gdy dane mają postać \emph{tabeli} --- a tak właśnie wyglądają cechy
połączeń sieciowych~\cite{grinsztajn2022why}. Sieci głębokie błyszczą przede
wszystkim tam, gdzie dane są surowe i mają wyraźną strukturę (obrazy, dźwięk,
tekst). Niniejsza praca sprawdza, czy ta obserwacja potwierdza się również w
naszym, ,,chmurowym'' zadaniu.
